\chapter{Úvod}

Cílem práce je vyvinout základní komponenty uživatelského rozhraní videohry žánru 2D plošinovka, postarat se o správnou funkčnost game manageru, který se stará například o přepínání mezi jednotlivými scénami, audio manageru, jehož úkolem je přehrávat všechny zvukové efekty a OST ve správný čas správným způsobem, a jednotlivých UI elementů.

Dále se zabývá tvorbou vizuálních efektů, které jsou v herním průmyslu důležité pro zkrášlení grafické stránky a přidání jistého realismu do zážitku z prostředí videohry. Vizuální efekty pomáhají hráčům chápat systémy, rychleji registrovat, co se děje, dostávat pozitivní/negativní zpětnou vazbu a vtažení do děje.

Úzce s vizuálními efekty souvisí i zvukové efekty -- snad každý vizuální efekt může mít svůj korespondující zvukový efekt, který ještě více potlačí účinek celé specifické akce. Audio je mezi průměrnými hráči často přehlíženo, protože se bere jako samozřejmost vzhledem k tomu, že zvukové odezvy se v reálném světě odehrávají všude kolem nás. Zkuste si ale někdy vaši oblíbenou videohru zahrát kompletně bez zvuku -- potěšení z videohry v takovém případě u mnoha titulů může klesnout drasticky. Práce se tedy zaměřuje jak na zvukové efekty, tak i na hudební složku.

Nakonec je nastíněn i jednoduchý model postprocessingu, který grafice naší hry dodává atmosféru, aby se tolik nepodobala retro videohrám jako je například herní titul Super Mario Bros \cite{super-mario-bros}.